\section{2025-01-14}


Worte können töten. Hetze gegen Minderheiten schafft ein Klima, in dem Gewalt legitim erscheint. Die ständige Wiederholung abwertender Phrasen erleichtert den Schritt vom Wort zur Tat und trug historisch zu Verfolgung und Völkermord bei. Auch im persönlichen Umfeld wirken Worte tödlich, wenn Beleidigungen, Mobbing und Demütigungen Menschen in Depression und Suizid treiben. Digitale Plattformen verstärken diese Gefahr, weil sich Hass und Falschmeldungen rasant verbreiten und Täter anonym bleiben können. Worte formen Denken und Handeln – jeder trägt Verantwortung, sie umsichtig zu wählen.
